% Sources: README "Numerical Parameters", "Regression Tests"; PLAN.md section 6

\section{Numerical scheme}
\label{app:numerics}

% OUTLINE: radial and vertical grids (nx, nz), launch-radius bisection, raycast;
% OUTLINE: C library vs Python mirror (<= 0.002 dex at the best fits);
% OUTLINE: regression tests.
% OUTLINE: test of the ray casting against the analytic projected area of an opaque finite
%          cone with self-shadowing (old manuscript, SEA_spectrum.zip, "Opaque finite cone with
%          self-shadowing"). The code tests only i <= theta_c (csrcs/test_wind_visibility.c);
%          the i > theta_c case needs a test in the code first.

\subsection{MCMC fitting}
\label{app:numerics:mcmc}
% Sources: ULXspectra-mc (branch mc-transfer-corrections, commit e0854b7): README "Adaptive Run Length", "Spectrum Thinning";
%          psrcs/likelihood.py, psrcs/params_mcmc.py (log_prior), psrcs/mcmc.py (walker start, best fit), psrcs/mcmc_adaptive.py
%          (last change of these files: d747a95). Run settings: run_settings.json of the 182 fits of the round mc_0p3-10keV
%          (ULXspectra-refit/results), read 2026-11-23.

\paragraph{Likelihood.}
The SEDs carry no measured errors (Sect.~\ref{sec:fitting:data}).
We therefore compare model and SED in $\log\nu L_\nu$ with one assumed scatter, $\sigma = 0.08$~dex, the same at all energies:
\begin{equation}
\ln\mathcal{L} = -\frac{1}{2}\sum_{k=1}^{n}\left[\left(\frac{\log y_k - \log m_k(\boldsymbol\theta)}{\sigma}\right)^2 + \ln\left(2\pi\sigma^2\right)\right],
\label{eq:loglike}
\end{equation}
% code: log_likelihood_with_blob() (psrcs/likelihood.py): chi = (log10 y_obs - log10 y_model)/sigma_log; sigma_log = 0.08 (ULX.conf DEFAULT)
where $y_k$ is the SED and $m_k$ the model, $\nu L_\nu$ at the energy $E_k$, and $\boldsymbol\theta$ the free parameters.
The first term is $\chi^2$.
The reduced $\chi^2$ that we quote is $\chi^2/(n - n_{\rm par})$ with this $\sigma$; it compares fits with each other but does not test them absolutely.
% code: red_chi2 = chi2/(n_points - n_params) (summaries/refits_summary.tsv; psrcs/plotting_mcmc.py)

\paragraph{Spectrum points.}
The SED files sample the spectrum much more densely in X-rays than at low energies.
The fit uses only the points nearest to energies evenly spaced in $\log E$, 20 per decade, so that every decade has the same weight in $\chi^2$; in 0.3--10~keV this gives $n = 30$ points.
% code: --thin logE, --points-per-decade 20 (README "Spectrum Thinning"); n_points = 30 for all fits of mc_0p3-10keV (refits_summary.tsv)
The SEDs are smooth model curves, not independent measurements: at 20 points per decade, linear interpolation reproduces them to 0.017~dex or better, well below $\sigma$.
More points would add no information on the shape and would only shrink the posterior widths as $1/\sqrt{n}$.
% source: README "Spectrum Thinning" ("Why 20": <= 0.017 dex; Fisher-matrix test, information per point within 5% (median) of
%         200 per decade)

\paragraph{Priors.}
The priors are uniform within the ranges of Table~\ref{tab:params} and zero outside them.
For a NS, a parameter set whose wind would be launched inside the NS ($\rin \le \RNS$) is also excluded.
% code: log_prior() (psrcs/params_mcmc.py): 0 inside the bounds, -inf outside; NS: launch_radius_cm(pars) <= R_NS -> -inf
With uniform priors and a constant $\sigma$, the maximum of the posterior is the minimum of $\chi^2$.

\paragraph{Sampler.}
We sample the posterior with the affine-invariant ensemble sampler \code{emcee} \citep{ForemanMackey2013} with 200 walkers.
Instead of the stretch move of \citet{Goodman2010}, which gets stuck between separate modes of the posterior, we use differential-evolution moves \citep{terBraak2006, terBraak2008}: 80\% of the proposals are $\boldsymbol\theta + \gamma(\boldsymbol\theta_j - \boldsymbol\theta_k)$, built from two other walkers $j$ and $k$, and 20\% are snooker moves.
% code: --moves de (README "Adaptive Run Length", "Further options"); nwalkers = 200 and moves = de in all 182 run_settings.json
The walkers start in a {small} Gaussian ball around {fixed starting values of the preset, which carry no physics and do not enter the prior}, with widths equal to the proposal scales of the parameters; a start point is kept only if its posterior is finite.
% code: psrcs/mcmc.py (p0 = theta0 + scale * N(0,1), accepted if log_prior and log_probability are finite)
{Since the whole ensemble starts in one small region, it can settle in the first mode it meets.
Nine NS fits that ended in a worse mode than an earlier fit were therefore re-run with half of the walkers around the earlier best fit and half spread uniformly over the prior.
For later fits the walkers start at the best of many points drawn uniformly over the prior, which gives a narrow mode near the edge of the prior a better chance to be found.}
% source: ULXspectra-mc docs/for_paper_session.md Note 13 (126dd71, edited 5c3a2b1): all rounds up to mc28_0p3-10keV start in the ball
%   around the ULX.conf starting values (e.g. log mdot0 = 1 +- 0.05); nine NS re-runs mc28b_0p3-10keV (100 walkers --init-from the
%   previous best, 100 --prior-walkers), running since 2026-11-24 13:14; later rounds: --prescan 5000, best 200 as starts
%   (code commits "Walkers start uniformly over the whole prior by default..." 848ff94 and "Pre-scan of the prior..." b4153cf).
% AUTHOR: after the mc28b re-runs finish, confirm that they replaced the nine mc28 fits; drop the last sentence if no later round is used.

\paragraph{Run length.}
Burn-in and run length are set by a convergence rule, checked every 100 steps.
The burn-in ends when it has lasted at least 300 steps, the median $\ln P$ of the walkers has changed by less than 1 since the previous check, and no walker lags behind.
A walker lags if its best $\ln P$ in the last 100 steps lies below the best of all walkers by more than $\chi^2_{0.9999}(n_{\rm par})/2$, i.e.\ far outside the posterior; during burn-in such walkers are moved onto copies of other walkers with a small jitter.
The reference is the best walker, not the median, because when a minority of walkers finds a much better mode, the median stays in the worse one and looks settled.
The production run ends when it is at least 20 times longer than the largest integrated autocorrelation time of the parameters, $\tau_{\rm max}$, and $\tau_{\rm max}$ has changed by less than 10\% since the previous check.
The effective sample size is then about $200 \times 20 = 4000$.
The production run stops after 5000 steps in any case, and the burn-in after 2500.
% code: psrcs/mcmc_adaptive.py; README "Adaptive Run Length" (burn_min 300, logp_tol 1, lag_prob 1e-4, n_tau 20, tau_rtol 0.1,
%       check_every 100, max_steps 5000 (production), burn_max = max_steps/2 (mcmc.py); rule recorded in run_settings.json "adaptive")
The best fit is the sample with the highest posterior in the whole production run.
% code: psrcs/mcmc.py (best = flat[argmax(flat_logp)])

In the 182 fits of 0.3--10~keV with the 117-run table of the transfer correction (Sect.~\ref{sec:results}), the burn-in took 300--1200 steps (median 400) and the production 700--5000 steps (median 1850).
% AUTHOR: these numbers are for the first round (v0); the calculation session is asked for those of the final round.
$\tau_{\rm max}$ was 12--570 steps (median 73), and the acceptance fraction 0.007--0.25 (median 0.10).
125 fits met the convergence rule; the other 57 stopped at 5000 production steps, with an effective sample size of 1600--3900.
% source: run_settings.json "adaptive" of the 182 fits mcmc_*.mc_0p3-10keV (ULXspectra-refit/results, read 2026-11-23):
%         converged 125/182; burn 300/400/1200, nsteps 700/1850/5000, tau_max 12.0/73.1/572.8, ess 1641/4061/7677 (unconverged: 1641-3867),
%         acceptance 0.0074/0.100/0.245 (min/median/max).
% AUTHOR: these numbers are for the current round (old table); to be updated after the C28 refits. Which of the 57
%         unconverged fits are best fits of Table bestfit_params (to be checked)?
